\section{Реализация алгоритмов}

\subsection{Быстрая сортировка на C++}

Реализация быстрой сортировки приведена в листинге~\ref{code:quicksort-cpp}.
Функция \texttt{partition} выбирает последний элемент подмассива в качестве опорного
(переменная \texttt{pivot}) и переставляет элементы так, чтобы все меньшие или равные
опорному оказались левее него, а большие "--- правее.
Функция \texttt{quickSort} рекурсивно применяет это разбиение.

\begin{listing}[H]
    \inputminted{cpp}{src/quicksort.cpp}
    \caption{Быстрая сортировка на языке C++}
    \label{code:quicksort-cpp}
\end{listing}

\subsection{Сортировка слиянием на Python}

Реализация сортировки слиянием на языке Python приведена в листинге~\ref{code:mergesort-py}.
Функция \mintinline{python}{merge_sort} рекурсивно делит список пополам до достижения
базового случая (список из одного элемента), после чего вспомогательная функция
\mintinline{python}{merge} последовательно сливает отсортированные части.

\begin{listing}[H]
    \inputminted{python}{src/mergesort.py}
    \caption{Сортировка слиянием на языке Python}
    \label{code:mergesort-py}
\end{listing}

\subsection{Ключевые отличия реализаций}

Реализация на C++ работает \textit{in-place}: функция \texttt{quickSort} изменяет
исходный вектор \texttt{arr} без выделения дополнительной памяти под временные массивы.
Реализация на Python является функциональной "--- каждый вызов \mintinline{python}{merge}
создаёт новый список, что обеспечивает чистоту кода, но требует $O(n)$ дополнительной памяти.
